% D.64. Independent root and PBC full proof and primary-source audits passed.
\subsection{An evaluated family with one additional ramification prime}

\begin{statement}{D.64}{Proved}{The one-primitive-prime Alexander calculation and its exact section fibres}
Let \(q\ne5\) be a rational prime of order twenty in
\((\mathbb Z/25)^\times\), and put \(S=\{5,q,\infty\}\).
For arbitrary \(a_5,a_q,b_5,b_q\in\mathbb Z_5\), set
\[
 a=a_5\kappa(5)+a_q\kappa(q),\qquad
 b=b_5\kappa(5)+b_q\kappa(q),\qquad s=(a,b).
                                                        \tag{D.64.1}
\]
Use the actual Alexander module \(A_s=\Lambda C\) of D.35,
with \(\mathfrak m=(5,u,v)\). Its unrestricted first
cohomology and its finite-ramification second cohomology are
\[
 \begin{gathered}
 H^0(G_{\mathbb Q},A_s)=0,\\
 H^1(G_{\mathbb Q},A_s)=H^1(G_{\mathbb Q,S},A_s)
                   \simeq\prod_{\mathbb N}\mathbb Z_5,\\
 H^2(G_{\mathbb Q,S},A_s)
       \xrightarrow[\mathrm{res}_q]{\ \sim\ }
 H^2(G_{\mathbb Q_q},A_s)=L_q.
 \end{gathered}                                         \tag{D.64.2}
\]
The first isomorphism of compact modules is noncanonical.
The second is the actual restriction map, with
\[
 L_q=A_s/\bigl((t_1^{a_q}t_2^{b_q}-1)A_s+
                    (q^{-1}\rho_s(\sigma_q)-1)A_s\bigr),
                                                        \tag{D.64.3}
\]
where the sum is closed and the Frobenius operator is
the one in D.35 and D.39. In the primitive axis case
\(b_q=0\), \(a_q\in\mathbb Z_5^\times\), this compact
module is itself a countable product of \(\mathbb Z_5\)'s.
The same conclusion holds after permuting the punctures.

Let \({\cal C}_S\) and \(D_s\) be the actual residue
module and operator of D.48. There is an algebraic splitting
\[
 H^2(G_{\mathbb Q},A_s)
       \simeq L_q\oplus({\cal C}_S/D_s{\cal C}_S).
                                                        \tag{D.64.4}
\]
It is specified by the inverse of restriction in (D.64.2),
followed by inflation to \(G_{\mathbb Q}\); the retraction
is localization at \(q\). In particular no closure is
inserted in the second summand. At each finite coefficient
level \(W_r=A_s/\mathfrak m^rA_s\), localization is an
isomorphism
\[
 H^2(G_{\mathbb Q},W_r)
       \xrightarrow{\ \sim\ }
            \bigoplus_{v\ne5}H^2(G_{\mathbb Q_v},W_r).
                                                        \tag{D.64.5}
\]
Consequently D.63's compact kernel \(\mathcal K_s\) is zero:
the separated unrestricted \(H^2\) is topologically exactly
the local family \(\mathcal L_s\) of (D.63.3). The Jannsen
kernel is exactly
\[
 \frac{{\cal C}_S\cap
             \prod_{v\notin S}(v^{-1}\rho_s(\sigma_v)-1)A_s}
      {D_s{\cal C}_S}.
                                                        \tag{D.64.6}
\]
It is nonzero, as in D.48.

For the geometrically pro-five metabelian quotient
\(M=F_5/F_5''\), the entire unrestricted section fibre
above (D.64.1) is nonempty if and only if
\[
                       a_qb_q(a_q-b_q)=0.
                                                        \tag{D.64.7}
\]
Every section in a nonempty fibre is represented over
\(G_{\mathbb Q,S}\), and that fibre is a torsor under
\(\prod_{\mathbb N}\mathbb Z_5\). Its global cusp classes
are exactly the three explicit packet maps in D.35,
with parameters in
\(\mathbb Z_5\kappa(5)+\mathbb Z_5\kappa(q)\).
Thus this is an evaluated four-parameter family within
Task D.5; arbitrary finite ramification sets and arbitrary
infinite ramification remain outside this assertion.
\end{statement}

\begin{proof}
We first determine the actual arithmetic group needed
for the calculation. This step also supplies the integral
lifting control that would be missing from a comparison
of rational ranks of finite augmentation quotients.

\paragraph{The maximal pro-five group and its single relation.}
Put \(K=\mathbb Q(\zeta_5)\), let \(G=G_{K,\{5,q\}}\),
and let \(P\) be its maximal pro-five quotient. Since
\(q\) has order four modulo five, there is just one prime
of \(K\) above \(q\); its residue cardinality is
\(Q=q^4\). Moreover
\[
                         v_5(Q-1)=1.                  \tag{D.64.8}
\]
The class-number, units, and Brauer calculation in D.1
gives
\[
 \dim_{\mathbb F_5}H^1(G_{K,\{5\}},\mathbb F_5)=3,
 \qquad H^2(G_{K,\{5\}},\mathbb F_5)=0.
\]
Here \(\mu_5\subset K\), so the residual Tate twist is
trivial. Finite-coefficient localization at the one new
prime, whose residue field contains \(\mu_5\), now gives
\[
       \dim H^1(G,\mathbb F_5)=4,
       \qquad \dim H^2(G,\mathbb F_5)=1.               \tag{D.64.9}
\]
The residue map in degree one is onto, and the degree-two
residue map is an isomorphism. The degree-three term
vanishes by odd-primary cohomological dimension two
\cite[I, Theorem 4.10(c); II, Proposition 2.9]{DE_MilneADT}.

For clarity, the bound on the pro-five relation rank does
not presuppose a comparison theorem for all coefficients.
If \(N=\ker(G\to P)\), then
\(H^1(N,\mathbb F_5)^P=0\). Otherwise the kernel of a
nonzero invariant homomorphism would give a quotient of
\(G\) which is an extension of \(P\) by \(\mathbb F_5\),
hence a pro-five quotient with nontrivial image of \(N\),
contrary to maximality. Hochschild--Serre therefore gives
\[
 H^1(P,\mathbb F_5)=H^1(G,\mathbb F_5),\qquad
 H^2(P,\mathbb F_5)\lhook\joinrel\longrightarrow
                                      H^2(G,\mathbb F_5).
                                                        \tag{D.64.10}
\]
Choose the images \(\tau,\sigma\) of a five-primary tame
inertia generator and an arithmetic Frobenius of \(K_q\).
They satisfy the actual relation
\[
                        \sigma\tau\sigma^{-1}=\tau^Q.
                                                        \tag{D.64.11}
\]
Their Frattini classes in \(P/\Phi(P)\) are independent.
Indeed the surjective residue map just noted supplies
a character nonzero on \(\tau\). The cyclotomic degree-five
extension \(K(\mu_{25})/K\) supplies a character zero on
\(\tau\) and nonzero on \(\sigma\), since its Frobenius
is \(q^4\not\equiv1\pmod{25}\). Thus they extend to four
minimal generators. Present \(P\) as a quotient of the
free pro-five group \(F=F(\sigma,\tau,h_1,h_2)\), with
kernel \(R\subset\Phi(F)\).

The element
\(r_0=\sigma\tau\sigma^{-1}\tau^{-Q}\) belongs to \(R\).
Its image in \(F^{\rm ab}\simeq\mathbb Z_5^4\) is
\((1-Q)e_\tau\), which has valuation exactly one.
On the other hand the image of
\(R^5[R,F]\), with closures understood, is contained
in \(25\mathbb Z_5^4\), because
\(R\subset\Phi(F)\) has abelianization contained in
\(5\mathbb Z_5^4\). Hence
\[
       0\ne[r_0]\in R/\overline{R^5[R,F]}.
                                                        \tag{D.64.12}
\]
The dual of this relation module is
\(H^2(P,\mathbb F_5)\): apply Hochschild--Serre to the
free presentation, using the minimality of its four
generators. By (D.64.9)--(D.64.10) the module is
one-dimensional. Thus \(r_0\) normally generates \(R\).
One can verify the last implication directly: a nonzero
quotient of \(R\) by its closed normal closure would be
detected in a finite five-group quotient of \(F\).
The resulting nontrivial normal subgroup has a nonzero
coinvariant Frattini quotient under that finite five-group;
the augmentation ideal of its group algebra is nilpotent.
This gives an \(F\)-invariant homomorphism
\(R\to\mathbb F_5\) vanishing on \(r_0\), contradicting
(D.64.12) and dimension one.

It follows that, as actual pro-five groups with the
specified local subgroup,
\[
       P=D*F(h_1,h_2),\qquad
       D=\langle\sigma,\tau\mid
                  \sigma\tau\sigma^{-1}=\tau^{q^4}\rangle
                \simeq\mathbb Z_5\rtimes_{q^4}\mathbb Z_5.
                                                        \tag{D.64.13}
\]
The last presentation has exactly the indicated
semidirect-product meaning. Collecting words puts them
in the form \(\tau^a\sigma^b\); the compact image of
\(\mathbb Z_5^2\) is closed and contains the dense word
subgroup. The homomorphism to the displayed semidirect
product distinguishes these normal forms. In particular
\(D\) is the actual maximal pro-five local group of
\(K_q\), and is embedded as a free factor of \(P\).

\paragraph{The coefficient comparison and the local degree-two isomorphism.}
For a finite five-primary \(P\)-module \(W\), restriction
for the free product gives
\[
                 H^n(P,W)\simeq H^n(D,W)\quad(n\ge2).
                                                        \tag{D.64.14}
\]
Here the other factor is free and has cohomological
dimension one. This is the finite-index-set case of
\cite[Theorem 4.2]{DE_GaertnerWingberg}; take the
corestricting subgroups to be the factors themselves.
Alternatively the universal property of crossed
homomorphisms decomposes the completed augmentation
ideal of a free product as the direct sum of the two
induced augmentation ideals. Dimension shifting and
Shapiro give precisely (D.64.14). The group \(D\) is
an extension of two copies of \(\mathbb Z_5\), so its
cohomological dimension is two. Thus the same is true
of \(P\).

For \(W=\mathbb F_5\), inflation \(P\to G\) is now an
isomorphism in every degree: degrees zero and one were
already identified, degree two follows from (D.64.9),
(D.64.10), and (D.64.12), and higher degrees vanish.
Every finite five-primary \(P\)-module has a filtration
with trivial \(\mathbb F_5\) factors. Its action factors
through a finite five-group, whose simple modules in
characteristic five are trivial. Induction on such a
filtration and the two cohomology long exact sequences
therefore prove inflation isomorphisms in all degrees
for every such \(W\). This proves the needed full/pro-five
comparison; it was not assumed in the relation-rank
calculation.

Each actual module \(W_r=A_s/\mathfrak m^rA_s\), restricted
to \(G\), is a \(P\)-module. Its augmentation and
five-adic filtrations have trivial residual factors
over \(K\), since \(\chi_5\equiv1\pmod5\); its finite
action consequently has five-power image. The same
observation for the graph of a crossed homomorphism
applies to all its finite cocycles. We may therefore
use (D.64.14) and the comparison just proved for these
actual twisted coefficients.

Over \(\mathbb Q\) the relevant quotient is an extension
\(\mathcal P\) of \(C_4=\operatorname{Gal}(K/\mathbb Q)\)
by \(P\). The local decomposition subgroup at \(q\)
surjects onto \(C_4\) and intersects \(P\) in the
embedded factor \(D\). Choose a complement \(T\simeq C_4\)
inside that local subgroup, by the coprime profinite
Schur--Zassenhaus theorem. In particular \(T\) normalizes
\(D\). The isomorphism (D.64.14) is \(T\)-equivariant.
Taking its invariants is exact because four is a unit
on every coefficient module. Local inflation from
\(D\) over \(K_q\) is also a cohomology isomorphism:
the discarded wild and prime-to-five tame factors have
zero positive five-primary cohomology, and the
remaining residue action \(Q\) lies in \(1+5\mathbb Z_5\).
Equivalently its two-generator tame resolution gives
the same finite cohomology. We obtain
\[
 H^2(G_{\mathbb Q,S},W_r)
      \xrightarrow[\mathrm{res}_q]{\ \sim\ }
                         H^2(G_{\mathbb Q_q},W_r).
                                                        \tag{D.64.15}
\]
All finite-ramification and all local finite-coefficient
cohomology groups are finite. Their towers are
Mittag--Leffler. The continuous comparison therefore
passes (D.64.15) to \(A_s\) without a derived-limit
correction. This proves the second line of (D.64.2),
with (D.64.3) supplied by the explicit local resolution
in D.63. The primitive-axis evaluation is D.43, since
\(q-1\) is a five-adic unit.

\paragraph{The integral rank of first cohomology.}
The zero-inertia argument of D.47, or D.48, identifies
unrestricted \(H^1\) with \(H^1(G_{\mathbb Q,S},A_s)\).
It is compact, countably based, and torsion-free by
D.35 and D.47. If \(a_q=b_q=0\), the action is already
unramified outside five, and D.36 proves the asserted
countably infinite rank.

Suppose now \((a_q,b_q)\ne(0,0)\). Applied over the finite
local extension \(K_q/\mathbb Q_q\), D.43 gives the
integral vanishings
\[
                       H^0(D,A_s)=H^1(D,A_s)=0.
                                                        \tag{D.64.16}
\]
Consequently every \(P\)-cocycle has a unique cohomologous
representative whose restriction to \(D\) is zero: subtract
the unique coboundary that is its restriction. By the
actual free product (D.64.13), the two values on
\(h_1,h_2\) can then be arbitrary elements of \(A_s\).
Indeed the universal property applies to their graphs
inside the pro-five semidirect product \(A_s\rtimes P\).
It gives a continuous bijection of compact modules,
hence a topological isomorphism
\[
              H^1(P,A_s)\simeq A_s\oplus A_s.
                                                        \tag{D.64.17}
\]
No quotient by residual coboundaries remains, by the
zero \(H^0\) in (D.64.16). Since \(T\) normalizes \(D\),
the space of these normalized cocycles is \(T\)-stable.
Exact averaging identifies the desired global group
with \(H^1(P,A_s)^T\).

Here is an integral argument that these invariants
have infinite rank. Reduce (D.64.17) modulo five and
filter it by powers of \(I=(u,v)\). The filtration is
\(T\)-stable: its terms are exactly the normalized
cocycles with all values in \(I^jA_s\), modulo five.
Let
\[
 W=H^1\bigl(P/\overline{\langle\!\langle D\rangle\!\rangle},
                                      \mathbb F_5\bigr).
\]
This is a two-dimensional \(T\)-module. On the associated
graded piece in degree \(j\), the coefficient action is
\(\mathbb F_5(j+2)^{j+1}\), and the \(P\)-action is
trivial. Evaluating the cocycle on words consequently
uses only the Frattini abelianization of the free
quotient. Thus the associated graded \(T\)-module is
\[
                 \mathbb F_5(j+2)^{j+1}\otimes W.
                                                        \tag{D.64.18}
\]
The four characters of \(T\) are defined over
\(\mathbb F_5\), and every representation is semisimple.
Choose a character occurring in the nonzero module
\(W\). For infinitely many \(j\), its product with
\(\chi_5^{j+2}\) is trivial. At each such \(j\), the
invariant dimension in (D.64.18) is at least \(j+1\).
Invariants are exact, and every finite augmentation
quotient of the actual module in (D.64.17) is a
surjective quotient. Therefore
\(H^1(P,A_s)^T/5\) surjects onto finite modules of
unbounded dimension. Its compact torsion-free
\(\mathbb Z_5\)-rank is infinite. D.47's countability
then identifies it with \(\prod_{\mathbb N}\mathbb Z_5\).
This argument uses actual arbitrary integral values on
the free generators; it does not infer inverse-limit
surjectivity from finite-index images at successive
augmentation levels.

\paragraph{The unrestricted obstruction and its separated kernel.}
At a fixed coefficient level, localization outside
\(S\) identifies the residue at \(v\notin S\) with
\(H^2(G_{\mathbb Q_v},W_r)\). If a global class is zero
locally at every \(v\ne5\), its residues outside \(S\)
vanish, so it comes from \(H^2(G_{\mathbb Q,S},W_r)\).
Its localization at \(q\) is zero, and (D.64.15) then
makes that class zero. Thus the localization map in
(D.64.5) is injective. It is surjective by the exact
Poitou--Tate argument in D.63: the component at five
can cancel the reciprocity functional, leaving every
prescribed finite-support family away from five.
This proves (D.64.5).

In particular \(H^2(G_{\mathbb Q,S},W_r)\) injects
into the unrestricted group, already because its
restriction at \(q\) is an isomorphism. The compact
term \({\cal L}_s\) in (D.48.2) therefore vanishes,
as does every finite kernel whose inverse limit is
\(\mathcal K_s\) in D.63. D.48 gives (D.64.6), and
D.63 supplies the asserted exact topological
identification of the separated quotient. The
restriction map at \(q\), together with (D.64.2),
is a retraction of the first term in (D.48.1).
Splitting that short exact sequence proves (D.64.4)
as an assertion about the actual, possibly
non-Hausdorff continuous cohomology groups.

\paragraph{Existence and the complete metabelian fibre.}
The arithmetic action on \(M\) and the abelian pair
both factor through \(G_{\mathbb Q,S}\). The pullback
of \(M\rtimes G_{\mathbb Q,S}\) over the exact abelian
section has kernel \(A_s\). It has a continuous
set-theoretic section: use the normal-form lift
\(x^{a(g)}y^{b(g)}\) and the arithmetic element \(g\).
Its multiplication defect is consequently an actual
continuous two-cocycle in \(A_s\). The extension has
a continuous group section if and only if its class
in \(H^2(G_{\mathbb Q,S},A_s)\) is zero. By (D.64.2),
this happens if and only if its local class at \(q\)
is zero.

Since \(q\not\equiv1\pmod5\), local Kummer theory gives
\(H^1(G_{\mathbb Q_q},\mathbb Z_5(1))=\mathbb Z_5\),
via valuation, and every local unit has zero completed
Kummer class. The local abelian pair is therefore
exactly \((a_q\kappa(q),b_q\kappa(q))\). If it is
nonzero, D.42 forces any metabelian lift into a cusp
packet and gives the necessity of (D.64.7).
Conversely each of its three lines is the abelian
image of the corresponding actual local cusp packet,
with an arbitrary completed tangent parameter.
The zero pair also has a unit cusp lift. Hence
(D.64.7) is exactly the local, and therefore the
global, existence criterion.

At every prime outside \(S\) an unrestricted section
above (D.64.1) has zero abelian inertia. D.42 makes
its entire metabelian inertia zero. Thus every such
section factors through \(G_{\mathbb Q,S}\). The
conjugacy and stabilizer analysis in D.35 identifies
its nonempty fibre with a torsor under the first
group in (D.64.2), without an additional abelian
stabilizer quotient. Finally the three full packet
formulas of D.35 show that a cusp section has its
parameter supported on \(\{5,q\}\) exactly when its
abelian pair does. They give precisely the cusp
classes asserted in the statement.
\end{proof}
